\subsection{实插值方法}
\label{sec:ch14-real-interpolation-method}
\setcounter{thmcount}{0}
\setcounter{equation}{0}

给定两个 Banach 空间 $B_0$ 和 $B_1$, 我们将定义在它们之间``插值''的 Banach 空间. 为简单起见, 假设 $B_1\subset B_0$. 例如, 可以取 $B_0=W_p^k(\Omega)$ 和 $B_1=W_p^m(\Omega)$, 其中 $k\leq m$. 对任意 $u\in B_0$ 和 $t>0$, 定义
\[
\MathBookFormulaID{formula-ch14-k-functional-definition}
K(t,u):=\inf_{v\in B_1}\bigl(\|u-v\|_{B_0}+t\|v\|_{B_1}\bigr).
\]
在某种意义下, $K$ 度量 $B_1$ 对 $u$ 的逼近程度, 其中逼近元的 $B_1$-范数乘以罚因子 $t$.

\MathBookSourcePage{380}
可以立即推出 $K(t,u)$ 的若干简单性质. 对 $u\in B_1$, 取 $v=u$ 得 $K(t,u)\leq t\|u\|_{B_1}$; 另一方面, 取 $v=0$ 得 $K(t,u)\leq\|u\|_{B_0}$. 对 $0<\theta<1$ 和 $1\leq p<\infty$, 定义范数
\MBSourceItem{1}
\begin{equation}
\MathBookFormulaID{formula-ch14-real-interpolation-norm}
\label{eq:ch14-real-interpolation-norm}
\|u\|_{[B_0,B_1]_{\theta,p}}
:=\left(\int_0^\infty t^{-\theta p}K(t,u)^p\,\frac{dt}{t}\right)^{1/p}.
\end{equation}
当 $p=\infty$ 时, 定义
\[
\MathBookFormulaID{formula-ch14-real-interpolation-infinity-norm}
\|u\|_{[B_0,B_1]_{\theta,\infty}}
:=\sup_{0<t<\infty}t^{-\theta}K(t,u).
\]
集合
\[
\MathBookFormulaID{formula-ch14-real-interpolation-space}
[B_0,B_1]_{\theta,p}=B_{\theta,p}
:=\{u\in B_0:\|u\|_{[B_0,B_1]_{\theta,p}}<\infty\}
\]
在范数~\eqref{eq:ch14-real-interpolation-norm} 下构成 Banach 空间(Bergh 与 Löfstrom 1976).

定义立即给出若干简单性质. 首先, 设 $A_i$ 和 $B_i$ 是上述两对 Banach 空间, 且 $B_i\subset A_i$($i=0,1$). 则 $B_{\theta,p}\subset A_{\theta,p}$(见习题~\ref{exr:ch14-01}). 其一个应用如下. 设 $\Omega\subset\widetilde\Omega$, 则
\MBSourceItem{2}
\begin{equation}
\MathBookFormulaID{formula-ch14-domain-restriction-interpolation-inclusion}
\label{eq:ch14-domain-restriction-interpolation-inclusion}
[W_p^k(\widetilde\Omega),W_p^m(\widetilde\Omega)]_{\theta,p}
\subset [W_p^k(\Omega),W_p^m(\Omega)]_{\theta,p},
\end{equation}
并有相应的范数不等式(见习题~\ref{exr:ch14-02}).

由于简化假设 $B_1\subset B_0$, 插值空间构成一个``尺度'': 总有 $B_1\subset B_{\theta,p}\subset B_0$(参见习题~\ref{exr:ch14-04}). 第二个包含关系部分地来自估计(参见习题~\ref{exr:ch14-03})
\MBSourceItem{3}
\begin{equation}
\MathBookFormulaID{formula-ch14-k-functional-interpolation-bound}
\label{eq:ch14-k-functional-interpolation-bound}
K(t,u)\leq C_{\theta,p}t^\theta\|u\|_{B_{\theta,p}}.
\end{equation}
这可以解释为, $u\in B_{\theta,p}$ 意味着 $u$ 在 $B_0$ 中能以 $t^\theta$ 阶被逼近. 事实上, 若 $p<\infty$, 则 $u\in B_{\theta,p}$ 还蕴含(习题~\ref{exr:ch14-05})
\MBSourceItem{4}
\begin{equation}
\MathBookFormulaID{formula-ch14-k-functional-little-o}
\label{eq:ch14-k-functional-little-o}
K(t,u)=o(t^\theta).
\end{equation}
这个更精细的逼近估计将在有限元应用一节中使用.

由~\eqref{eq:ch14-k-functional-interpolation-bound} 立即可知, 对所有 $1\leq p\leq\infty$, 有 $B_{\theta,p}\subset B_{\theta,\infty}$. 此外, 已知对所有 $1\leq p\leq\infty$, 有 $B_{\theta,1}\subset B_{\theta,p}$(Bergh 与 Löfstrom 1976). 若 $\theta_1\leq\theta_2$, 则对所有 $1\leq p\leq\infty$, 有 $B_{\theta_2,p}\subset B_{\theta_1,p}$; 若 $p\leq q$, 则对所有 $0<\theta<1$, 有 $B_{\theta,p}\subset B_{\theta,q}$(Bergh 与 Löfstrom 1976).
插值理论的关键结果涉及 Banach 空间上的算子.

\MathBookSourcePage{381}
\MBSourceItem{5}
\begin{Proposition}
\label{prop:ch14-operator-interpolation}
设 $A_i$ 和 $B_i$ 是上述两对 Banach 空间, 且线性算子 $T$ 将 $A_i$ 映到 $B_i$($i=0,1$). 则 $T$ 将 $A_{\theta,p}$ 映到 $B_{\theta,p}$. 而且
\MBSourceItem{6}
\begin{equation}
\MathBookFormulaID{formula-ch14-operator-interpolation-bound}
\label{eq:ch14-operator-interpolation-bound}
\|T\|_{A_{\theta,p}\to B_{\theta,p}}
\leq \|T\|_{A_0\to B_0}^{1-\theta}\|T\|_{A_1\to B_1}^{\theta}.
\end{equation}
\end{Proposition}

\begin{Proof}
令 $M_i:=\|T\|_{A_i\to B_i}$. 对任意 $v\in A_1$,
\[
\MathBookFormulaID{formula-ch14-operator-k-functional-bound}
\begin{aligned}
K_B(t,Tu)&\leq\|Tu-Tv\|_{B_0}+t\|Tv\|_{B_1}\\
&\leq M_0\|u-v\|_{A_0}+tM_1\|v\|_{A_1}.
\end{aligned}
\]
对 $v\in A_1$ 取下确界, 得
\[
\MathBookFormulaID{formula-ch14-scaled-k-functional-bound}
K_B(t,Tu)\leq M_0K_A(tM_1/M_0,u).
\]
对此积分, 并在不等式右端的积分中作变量替换 $s=tM_1/M_0$, 即得结论.
\end{Proof}

还有两个关于实插值方法的结果是我们所需但不予证明的. 证明见 Bergh 与 Löfstrom (1976). 第一个称为``重申定理'', 它断言
\MBSourceItem{7}
\begin{equation}
\MathBookFormulaID{formula-ch14-reiteration-theorem}
\label{eq:ch14-reiteration-theorem}
[[B_0,B_1]_{\theta_0,p_0},[B_0,B_1]_{\theta_1,p_1}]_{\lambda,q}
=[B_0,B_1]_{(1-\lambda)\theta_0+\lambda\theta_1,q}
\end{equation}
对任意 $0\leq\theta_0<\theta_1\leq1$, $1\leq p_0,p_1,q\leq\infty$ 和 $0<\lambda<1$ 成立. 第二个结果涉及对偶空间. 假设 $B_1$ 在 $B_0$ 中稠密, 则
\MBSourceItem{8}
\begin{equation}
\MathBookFormulaID{formula-ch14-interpolation-duality}
\label{eq:ch14-interpolation-duality}
[B_0,B_1]_{\theta,p}'=[B_1',B_0']_{1-\theta,p'}
\end{equation}
对任意 $0<\theta<1$ 和 $1\leq p<\infty$ 成立, 其中 $1/p+1/p'=1$. 下一节将把这些结果用于负指标 Sobolev 空间.

最后指出, 限制 $B_1\subset B_0$ 并非必要. 一般只需 $v_0+v_1$ 对 $v_i\in B_i$ 有定义. 此时
$[B_0,B_1]_{\theta,p}=[B_1,B_0]_{1-\theta,p}$. 因而对偶定理也可写为
$[B_0,B_1]_{\theta,p}'=[B_0',B_1']_{\theta,p'}$.
